\documentclass[11pt]{article}

\usepackage{amsmath}
\usepackage{amssymb}
\usepackage{amsfonts}
\usepackage{hyperref}
\usepackage{geometry}
\usepackage{cite}

\geometry{margin=1in}

\title{Wave-Packet Propagation and Constraint Stability on a Frozen Cochain Operator Architecture}
\author{Tomislav Rupi\'c\\Independent Researcher\\\v{S}ibenik, Croatia}
\date{March 2026}

\begin{document}
\maketitle

\begin{abstract}
We document Stage 9 wave-dynamics diagnostics on the previously validated kernel-induced cochain architecture. The study is limited to second-order packet evolution in the frozen scalar sector and the frozen projected transverse sector; first-order Dirac--Kaehler packet dynamics is analyzed separately in Stage 9B. Localized packets are evolved with the existing operators $L_0$ and $L_1$, and the diagnostics record packet center motion, width evolution, norm and energy behavior, anisotropy, divergence residuals, and redundancy under exact longitudinal contamination. Across periodic domains of size $n=16,20,24$, the scalar sector shows coherent packet propagation with controlled spreading and relative energy drift at the level of $10^{-3}$. The projected transverse sector shows similarly coherent propagation while preserving the divergence constraint and transverse projection to numerical precision. A separate redundancy test confirms that longitudinally shifted initial data evolve identically after projection, with trajectory, width, and energy differences at machine precision. These results remain strictly operator-level: they document reproducible transport and constraint stability on a frozen discrete architecture, but do not establish particles, gauge symmetry, relativistic invariance, or continuum field dynamics.
\end{abstract}

\section{Introduction}
The HAOS/IIP numerical program has so far established a layered operator architecture on a common Gaussian interaction substrate. Earlier stages validated a kernel-induced scalar branch with continuum Laplacian scaling, a restricted transverse edge sector with a stable low spectral band, a set of local and nonlocal constraint diagnostics, and a consistent Dirac--Kaehler factorization on the same cochain complex. Those results were algebraic, spectral, or static in character. They established what operators exist on the frozen architecture and which constraint relations remain numerically stable, but they did not yet test whether the same architecture supports coherent dynamical transport of localized disturbances.

Stage 9 addresses that gap. The objective is deliberately narrow: determine whether the previously validated scalar and projected transverse sectors admit wave-packet evolution that remains structured over a short propagation window. The focus is on behavioral morphology rather than physical interpretation. In particular, the numerical tests ask whether localized packets translate or spread in a controlled way, whether the projected transverse sector remains closed under evolution, and whether exact longitudinal contamination is removed consistently by the same projection mechanism already used in the static diagnostics.

This paper therefore does not attempt reconstruction of continuum physics, identification of particles, or proof of gauge symmetry. Its purpose is to document a reproducible dynamical benchmark on a frozen architecture. The question is not what the system means, but what the already validated operators do when they are evolved in time. Stage 9 investigates second-order transport in scalar and projected transverse sectors only; first-order Dirac--Kaehler packet dynamics is analyzed separately in Stage 9B.

\section{Numerical Architecture}
No operator modification was introduced for Stage 9. All dynamics are built from the previously validated periodic cochain architecture generated by the Gaussian kernel substrate. The scalar sector uses the frozen node operator $L_0$, while the edge sector uses the frozen Hodge operator
\begin{equation}
L_1 = d_0 d_0^{*} + d_1^{*} d_1.
\end{equation}
The transverse runs are restricted to the same projected subspace used in the earlier static diagnostics,
\begin{equation}
\ker(d_0^{*}) \cap (\mathcal{H}^1)^\perp,
\end{equation}
with exact and harmonic components removed through the established projector.

The Stage 9 experiments were run on periodic domains with resolutions
\begin{equation}
n = 16,\ 20,\ 24,
\end{equation}
using the frozen kernel parameter $\epsilon = 0.2$, packet width $\sigma = 0.08$, packet center $(0.25,0.5,0.5)$, and an initial kick directed along the $x$ axis. The time-evolution window was fixed at
\begin{equation}
t_{\mathrm{final}} = 0.45,
\end{equation}
and the time step was chosen adaptively from the spectral radius estimate
\begin{equation}
\Delta t = 0.45\,\frac{2}{\sqrt{\lambda_{\max}}}.
\end{equation}
In the committed runs this choice produced two explicit update steps per resolution.

Both sectors were evolved with the same second-order scheme,
\begin{equation}
\partial_t^2 q = -L q,
\end{equation}
implemented in velocity-Verlet form. For the scalar sector, $q = \phi$ and $L=L_0$. For the projected edge sector, $q=A$ and $L=L_1$, with projection applied to the field, velocity, and acceleration after each update so that the trajectory remains in the restricted transverse subspace. Packet observables were evaluated from the squared amplitude profile $|q|^2$ and included center trajectory, effective width, Euclidean norm, wave-energy surrogate,
\begin{equation}
E(q,v) = \frac{1}{2}\big(\langle v,v\rangle + \langle q, L q\rangle\big),
\end{equation}
anisotropy ratio, and, for the edge sector, the divergence residual $\|d_0^{*}A(t)\|$ and projection residual $\|P_{\perp}A(t)-A(t)\|$.

Throughout this paper the discrete norm is the Euclidean coefficient norm
\begin{equation}
\|q\| = \sqrt{\langle q,q\rangle},
\end{equation}
with the same coefficient inner product used in the numerical implementation. The second-order energy functional is the quadratic form
\begin{equation}
E(q,v) = \frac{1}{2}\big(\langle v,v\rangle + \langle q, L q\rangle\big),
\end{equation}
evaluated with the frozen operator of the sector under study. In the absence of integration error this is the natural conserved quantity for the continuous second-order evolution generated by the self-adjoint operator $L$, so Stage 9 reports energy drift rather than raw norm drift as the primary stability diagnostic.

Table~\ref{tab:common-parameters} summarizes the common numerical settings. These settings were fixed across the Stage 9A diagnostics so that the comparison remains behavioral rather than architectural.

\begin{table}[t]
\centering
\begin{tabular}{ll}
\hline
Quantity & Value \\
\hline
Domain type & periodic graph-based cochain complex \\
Resolutions & $n = 16,20,24$ \\
Kernel width & $\epsilon = 0.2$ \\
Packet center & $(0.25,0.5,0.5)$ \\
Packet width & $\sigma = 0.08$ \\
Kick axis & $x$ \\
Kick strength & $1.0$ \\
Final time & $t_{\mathrm{final}} = 0.45$ \\
Time-step rule & $0.45\,2/\sqrt{\lambda_{\max}}$ \\
Integrator & velocity-Verlet / leapfrog form \\
\hline
\end{tabular}
\caption{Common Stage 9 parameters used in the scalar, transverse, and constraint-stability diagnostics.}
\label{tab:common-parameters}
\end{table}

\section{Scalar Wave-Packet Diagnostics}
The scalar diagnostic starts from a localized Gaussian profile on the node sector,
\begin{equation}
\phi(x,0) \propto \exp\!\left(-\frac{|x-x_0|^2}{2\sigma^2}\right),
\end{equation}
with a normalized velocity profile aligned with the $x$ direction. The evolution is then generated directly by the frozen scalar operator through
\begin{equation}
\partial_t^2 \phi = -L_0\phi.
\end{equation}
The purpose of this block is modest: determine whether the packet translates and broadens smoothly, rather than fragmenting or becoming numerically unstable.

Across all three tested resolutions, the scalar packet shows a clear positive displacement of its center along the kick axis. The total center shift over the recorded time window is
\begin{equation}
0.05197,\ 0.05006,\ 0.04903
\end{equation}
for $n=16,20,24$, respectively. The width change is small compared with the initial packet size, ranging from $+9.95\times 10^{-4}$ at $n=16$ to $-1.74\times 10^{-3}$ at $n=24$. The energy drift remains at the level of
\begin{equation}
1.23\times 10^{-3},\ 1.46\times 10^{-3},\ 1.12\times 10^{-3},
\end{equation}
while the maximal anisotropy ratio stays between $1.38$ and $1.44$.

Over the tested short-time window, packet center motion is approximately linear; no asymptotic transport regime is inferred. Anisotropy is monitored here as a morphology diagnostic; no refinement or continuum-limit claim is made.

The Euclidean field norm is not conserved in this formulation and changes by a few percent over the run, increasing from about $2.2\%$ at $n=16$ to about $7.8\%$ at $n=24$. For that reason the energy surrogate, rather than the raw $L^2$ norm, is the more reliable stability indicator in this block. On that metric the evolution is well controlled. Combined with the smooth center motion and bounded width variation, the scalar run supports a simple behavioral reading: over the tested short-time window the frozen scalar sector admits coherent packet transport with mild dispersive reshaping and no visible fragmentation.

\section{Transverse Sector Propagation}
The transverse diagnostic uses the same packet geometry but places the disturbance in the projected edge sector. A localized edge profile is first built on oriented edge midpoints and then projected into the restricted transverse subspace. The evolution equation is
\begin{equation}
\partial_t^2 A = -L_1 A,
\end{equation}
with the transverse projector enforced at each substep. The monitored observables are the same transport measures used in the scalar case, augmented by the divergence residual and projection residual.

The committed Stage 9 results show coherent motion of the packet center and controlled width evolution across all three resolutions. The center shift in the $x$ direction is
\begin{equation}
0.04718,\ 0.04526,\ 0.04414
\end{equation}
for $n=16,20,24$. The width narrows modestly over the same window, with changes of
\begin{equation}
-1.28\times 10^{-3},\ -4.27\times 10^{-3},\ -5.99\times 10^{-3}.
\end{equation}
Relative energy drift remains small,
\begin{equation}
1.06\times 10^{-3},\ 1.39\times 10^{-3},\ 1.09\times 10^{-3},
\end{equation}
and the maximal anisotropy ratio stays close to $1.5$.

Over the tested short-time window, packet center motion is approximately linear; no asymptotic transport regime is inferred. Anisotropy is monitored here as a morphology diagnostic; no refinement or continuum-limit claim is made.

The decisive metric in this section is not transport alone, but closure of the projected sector under dynamics. The divergence residual
\begin{equation}
\max_t \|d_0^{*}A(t)\|
\end{equation}
remains between $1.44\times 10^{-16}$ and $1.74\times 10^{-16}$, while the maximal projection residual
\begin{equation}
\max_t \|P_{\perp}A(t)-A(t)\|
\end{equation}
remains between $1.48\times 10^{-16}$ and $6.96\times 10^{-16}$. These numbers are at numerical roundoff. In practical terms, the packet does not acquire measurable longitudinal contamination under the frozen projected dynamics.

This makes the transverse result stronger than a generic edge-wave observation. The packet not only propagates coherently, but does so while remaining inside the restricted transverse architecture already validated in the static stages. That is the narrow Stage 9 claim for this block: the projected transverse sector is dynamically viable on the frozen architecture, and its defining divergence constraint remains closed under the tested evolution.

\section{Constraint Stability and Redundancy Test}
The third Stage 9 block tests whether the exact-edge redundancy documented in the earlier static diagnostics remains stable during time evolution. Starting from the same projected transverse packet, an exact component is added to the initial data,
\begin{equation}
A_{\mathrm{cont}}(0) = A_{\mathrm{proj}}(0) + d_0\varphi,
\end{equation}
with a matching perturbation added to the initial velocity. Both the baseline and contaminated initial states are then projected and evolved under the same transverse dynamics.

The comparison is intentionally strict. Rather than checking only a spectral quantity, the run measures trajectory differences, width differences, and energy differences between the two projected evolutions. The result is numerically exact for all practical purposes. Across the three tested resolutions, the maximal center difference never exceeds
\begin{equation}
4.44\times 10^{-16},
\end{equation}
the maximal width difference never exceeds
\begin{equation}
1.11\times 10^{-16},
\end{equation}
and the maximal energy difference never exceeds
\begin{equation}
6.66\times 10^{-16}.
\end{equation}
The divergence and projection residuals in the baseline trajectory remain at the same roundoff level reported in the previous section.

These numbers support a conservative but meaningful interpretation. The frozen projection mechanism does not merely clean the initial data once; it closes the dynamics under the equivalence generated by exact contamination. In the language used throughout the repository, the longitudinal redundancy remains dynamically consistent. In less interpretive terms, projected evolution depends only on the restricted transverse content of the initial data and is insensitive, to machine precision, to the added exact component.

\section{Behavioral Summary}
Taken together, the three Stage 9 diagnostics support a clear behavioral classification of the tested short-time regime. The scalar packet exhibits directed center motion with small energy drift and mild reshaping. The projected transverse packet shows comparable center motion with equally small energy drift, but with the additional property that its divergence residual and projection residual remain at numerical precision throughout the run. The separate redundancy test confirms that exact longitudinal contamination is removed consistently before and during propagation.

The most restrained summary is therefore morphological rather than interpretive. Over the tested time window, the frozen architecture supports coherent directed transport with dispersive spreading, stable constraint preservation, and robust sector decomposition. The runs are short, the transport has not yet been mapped over a broad spectral range, and no asymptotic transport law is asserted here. What can be said is narrower and sufficient: the packets move in the kicked direction, remain organized, and do not collapse into noise or leak out of the constrained subspace.

A second important point is that the scalar and transverse sectors differ in what must remain stable. In the scalar case, the key diagnostic is smooth propagation with bounded energy drift. In the transverse case, the key diagnostic is smooth propagation while preserving the projected constraint structure. The Stage 9 run succeeds on both counts.

\section{Interpretation Boundary}
Stage 9 is an operator-level dynamics study only. It does not establish particle ontology, photon-like excitations, fermion sectors, gauge-symmetry proofs, relativistic invariance, or emergent spacetime reconstruction. The committed diagnostics show only that localized disturbances evolved by the frozen scalar and projected transverse operators propagate coherently over the tested window, and that the transverse constraint structure remains numerically closed under that evolution.

This boundary matters because the architecture was intentionally frozen before the dynamics study began. The purpose of Stage 9 was to test what the existing operators do, not to redesign them until a preferred interpretation appears. The resulting claims must therefore remain at the level of transport, spreading, constraint preservation, and redundancy consistency.

\section{Next Diagnostic Direction}
The natural continuation of Stage 9 is not immediate interpretation, but broader behavioral classification. Four next diagnostics follow directly from the current results. First, multi-packet runs can test whether the same architecture supports merging, scattering, trapping, or interference regimes beyond single-packet motion. Second, spectral regime mapping can connect packet morphology to the low and intermediate bands already identified in the static studies. Third, perturbation tests can determine how transport and constraint stability respond to mild anisotropy, graph disorder, or kernel variation without modifying the core operators. Fourth, the already validated Dirac--Kaehler sector can be added as a separate first-order dynamics block, provided it is held to the same standard of bounded interpretation.

In repository terms, these directions belong to a behavioral atlas rather than a new architecture stage. Stage 9 establishes that packet transport is worth classifying; the next stages should map when that transport remains coherent, when it becomes dispersive, and when the existing constraint structure begins to degrade.

\section{Conclusion}
Stage 9 confirms that the frozen HAOS/IIP architecture supports coherent transport regimes with stable constraint structure. On periodic domains of size $n=16,20,24$, scalar packets evolved by the frozen operator $L_0$ move smoothly with bounded energy drift and controlled spreading. Projected transverse packets evolved by the frozen edge operator $L_1$ show similarly coherent transport while maintaining divergence and projection residuals at numerical precision. A separate redundancy test demonstrates that longitudinally shifted initial data evolve identically after projection.

The conclusion is intentionally narrow. These are behavioral properties of a discrete interaction architecture, not claims about particles or continuum field theories. What Stage 9 establishes is that the previously validated static structure also supports reproducible wave-packet transport and internal constraint closure under explicit time evolution. That is enough to justify the next step: systematic mapping of the pre-geometric dynamical landscape on the same frozen architecture.

\appendix
\section{Resolution Summary}
Table~\ref{tab:resolution-summary} collects the main Stage 9A scalar and transverse metrics used in the body text.

\begin{table}[ht]
\centering
\begin{tabular}{cccccc}
\hline
Sector & $n$ & center shift & width change & rel. energy drift & max constraint \\
\hline
scalar & 16 & 0.05197 & $+9.95\times10^{-4}$ & $1.23\times10^{-3}$ & -- \\
scalar & 20 & 0.05006 & $-7.61\times10^{-4}$ & $1.46\times10^{-3}$ & -- \\
scalar & 24 & 0.04903 & $-1.74\times10^{-3}$ & $1.12\times10^{-3}$ & -- \\
transverse & 16 & 0.04718 & $-1.28\times10^{-3}$ & $1.06\times10^{-3}$ & $1.61\times10^{-16}$ \\
transverse & 20 & 0.04526 & $-4.27\times10^{-3}$ & $1.39\times10^{-3}$ & $1.74\times10^{-16}$ \\
transverse & 24 & 0.04414 & $-5.99\times10^{-3}$ & $1.09\times10^{-3}$ & $1.44\times10^{-16}$ \\
\hline
\end{tabular}
\caption{Main Stage 9A transport and constraint metrics. The constraint column records $\max_t\|d_0^{*}A(t)\|$ for the transverse sector.}
\label{tab:resolution-summary}
\end{table}

\begin{thebibliography}{9}

\bibitem{paper1}
Tomislav Rupi\'c,
\newblock \emph{Kernel-Induced Operator Architecture and the Emergence of a Transverse Spectral Band on Discrete Substrates},
\newblock HAOS-IIP preprint, 2026.

\bibitem{paper2}
Tomislav Rupi\'c,
\newblock \emph{Continuum Reconstruction and Scalar-Transverse Coupling in Kernel-Induced Operator Systems},
\newblock HAOS-IIP preprint, 2026.

\bibitem{paper3}
Tomislav Rupi\'c,
\newblock \emph{Dirac-Kaehler Factorization on a Validated Cochain Operator Architecture},
\newblock HAOS-IIP preprint, 2026.

\bibitem{paper4}
Tomislav Rupi\'c,
\newblock \emph{Gauge-Like Diagnostics on a Kernel-Induced Cochain Operator Architecture},
\newblock HAOS-IIP preprint, 2026.

\bibitem{repo}
Tomislav Rupi\'c,
\newblock \emph{HAOS-IIP Repository and Stamped Numerical Artifacts},
\newblock GitHub repository, \url{https://github.com/tomislavrupic}.

\end{thebibliography}

\end{document}
